\newglossaryentry{bibliotheque} {
  name={bibliothèque},
  type=definitions,
  description={Dans le contexte du séquençage d'ADN, désigne un ensemble de fragments
  d'ADN préparés en vue du séquençage (\textit{library})}
}

\newglossaryentry{cenancetre} {
  name={cénancêtre},
  type=definitions,
  description={Espèce la plus récente que deux taxons ont en commun. Aussi appelé dernier
  ancêtre common ou DAC (\textit{last common ancestor, LCA})}
}

\newglossaryentry{contig} {
  name={contig},
  type=definitions,
  description={Séquence d'ADN continue assemblée à partir de fragments séquencés séparément qui appartenaient à la même séquence avant le processus de séquençage}
}

\newglossaryentry{hyperparametre} {
  name={hyperparamètre},
  type=definitions,
  description={Valeur numérique configurée avant l'entraînement d'un modèle d'apprentissage automatique qui détermine comment ses paramètres sont ajustés durant l'entraînement}
}

\newglossaryentry{indel} {
  name={indel},
  type=definitions,
  description={Contraction de insertion / délétion, un type de modification de l'ADN qui consiste à substituer ou éliminer une ou plusieurs nucléotides dans une séquence}
}

\newglossaryentry{lecture} {
  name={lecture},
  type=definitions,
  description={Dans le contexte du séquençage d'ADN, désigne une séquence brute
      d'un fragment d'ADN obtenue avec une technologie de séquençage (\textit{read})}
}

\newglossaryentry{metagenome} {
  name={métagénome},
  type=definitions,
  description={Ensemble du matériel génétique d'un microbiome \cite{roy2024}}
}

\newglossaryentry{metagenomique} {
  name={métagénomique},
  type=definitions,
  description={Étude des métagénomes \cite{roy2024}}
}

\newglossaryentry{microbiome} {
  name={microbiome},
  type=definitions,
  description={Ensemble de microorganismes vivants dans un habitat donné ainsi que leurs activités et propriétés physio-chimiques (e.g. métabolites, éléments génétiques mobiles, virus, etc.) \cite{Berg2020}}
}

\newglossaryentry{microbiote} {
  name={microbiote},
  type=definitions,
  description={Ensemble de microorganismes vivants dans un habitat donné (e.g. bactéries, archées, protistes, fonges, algues, etc.) \cite{Berg2020}}
}

\newglossaryentry{parametre} {
  name={paramètre},
  type=definitions,
  description={Valeur numérique contenue dans un modèle d'apprentissage automatique ajustée durant l'entraînement qui détermine le comportement du modèle lors de l'inférence}
}
